% !TEX program = xelatex
% 编译方式：把本文件夹整包复制到纯 ASCII 路径后，执行两遍 xelatex（详见 docs/_manual_common/README.md）。
% 首版骨架：内容取自 src/analysis/ 与 include/hacdcpf/analysis/ 源码
%           （hosting_capacity / multidimensional_weak_link / counterfactual_planning），公式与参数以代码为准。
\documentclass[UTF8,fontset=windows,a4paper,11pt]{ctexart}
\usepackage{hysim_manual}

\title{\textbf{交直流混合高弹性能源电力系统仿真平台}\\[0.5em]
       {\Large 承载力与薄弱环节分析技术手册}\\[0.3em]
       {\large 数学模型、接口参数与验证校核（首版骨架）}}
\author{HySim-XJTU-HRPES（西安交通大学 HRPES 团队）}
\date{\today}

\begin{document}
\maketitle
\begin{abstract}
\noindent 本手册依据 HySim-XJTU-HRPES 平台分析模块
（\srcpath{src/analysis/}、\srcpath{include/hacdcpf/analysis/}）整理，覆盖三类面向
规划与评估的后处理方法：依据《DL/T 2041-2025》的设备级承载力评估、多维度薄弱环节
筛查、反事实规划措施评估。全部结论以源码为准：承载力校核默认关闭（仅做规划筛查），
薄弱环节为非补偿式多准则筛查与决策路由（不宣称归因量等于反事实投资收益、缺失维度不
补零），反事实措施施加于系统副本、经济/可靠性/弹性指标使用选项给定的有界预算与时窗。
命名空间为 \texttt{hacdcpf::analysis}。文档版式与《碳流追踪与碳排放核算技术手册》一致。
\end{abstract}

\tableofcontents
\newpage

\section{阅读指南与约定}
\subsection{数据来源}
\begin{itemize}
  \item \srcpath{include/hacdcpf/analysis/hosting\_capacity.hpp}、
        \srcpath{src/analysis/hosting\_capacity.cpp} —— 承载力评估接口与实现；
  \item \srcpath{include/hacdcpf/analysis/multidimensional\_weak\_link.hpp}、
        \srcpath{src/analysis/multidimensional\_weak\_link.cpp} —— 多维薄弱环节筛查；
  \item \srcpath{include/hacdcpf/analysis/counterfactual\_planning.hpp}、
        \srcpath{src/analysis/counterfactual\_planning.cpp} —— 反事实规划措施评估；
  \item \srcpath{docs/capacity\_analysis\_implementation.md}、
        \srcpath{docs/multidimensional\_weak\_link\_identification.md} —— 既有契约文档。
\end{itemize}
以上头文件与 \srcpath{include/hacdcpf/analysis/} 同目录，其余同名头文件
（\srcpath{short\_circuit.hpp}、\srcpath{scenario\_generation.hpp} 等）实现于兄弟模块，
不属于本手册范围。

\subsection{单位制与索引空间}
承载力 $S_d$ 单位 MVA/MW，功率因数 $\cos\theta$ 无量纲；薄弱环节压力量归一化后无量纲，
时间维度以选项声明的评估周期为索引空间；反事实四维收益向量的物理量分别为经济运行成本
（\$/yr）、碳排放（tCO\textsubscript{2}/yr）、可靠性 EENS（MWh/yr）与弹性加权失负荷 ENS（MWh）。
承载力评估以变压器下游供电区（BFS）为索引，结果按 \fld{transformers} 与 \fld{areas} 分列。

\subsection{诚实口径}
三类分析均为规划/评估后处理，不改变调度：承载力 \fld{enable\_verification} 默认为
\texttt{false}，即默认只给规划筛查值、不回跑潮流/短路/谐波引擎校核；薄弱环节为非补偿式
筛查，缺失维度不被当作零、也不宣称归因排放/失负荷等于反事实投资收益；反事实措施施加于
系统副本，经济/可靠性/弹性指标使用 \fld{reliability\_physical\_budget}、
\fld{resilience\_horizon\_hours} 等有界预算与时窗。相关限制在 \S 5.2 如实列出。

\section{功能定位与架构}
\subsection{公共入口族}
\begin{center}\footnotesize
\begin{tabular}{@{}L{6.6cm}L{3.0cm}L{5.4cm}@{}}
\toprule
入口 & 定义位置 & 说明\\
\midrule
\srcpath{assess\_hosting\_capacity(sys, opt)} & hosting\_capacity.hpp & 依《DL/T 2041-2025》做设备级承载力评估（可选校核）\\
\srcpath{hosting\_capacity\_result\_to\_json(r)} & 同上 & 承载力结果 JSON 序列化\\
\srcpath{run\_multidimensional\_weak\_link\_assessment(...)} & multidimensional\_weak\_link.hpp & 四维薄弱环节筛查（entities/periods/options）\\
\srcpath{generate\_counterfactual\_measures(system, options)} & counterfactual\_planning.hpp & 候选措施生成（线路/储能/联络开关/自动化/DER）\\
\srcpath{apply\_counterfactual\_measure(system, measure)} & 同上 & 将单条措施施加到系统副本\\
\srcpath{run\_counterfactual\_planning\_assessment(...)} & 同上 & 基线 + 逐措施 + 两两协同评估\\
\bottomrule
\end{tabular}
\end{center}

\subsection{关键数据结构}
\compmeta{HostingCapacityResult}{include/hacdcpf/analysis/hosting\_capacity.hpp}
主要字段：\fld{standard}（默认 \texttt{"DL/T 2041-2025"}）、\fld{transformers}
（\texttt{TransformerHostingResult} 数组）、\fld{areas}（\texttt{AreaHostingResult} 数组）、
\fld{warnings}、\fld{verification}（\texttt{HostingVerification}）。

\compmeta{MultidimensionalWeakLinkResult}{include/hacdcpf/analysis/multidimensional\_weak\_link.hpp}
主要字段：\fld{schema}、\fld{dimension\_coverage}（四维覆盖）、\fld{entities}、
\fld{critical\_periods}、\fld{pareto\_count}、\fld{compatible\_count}、\fld{conflict\_count}。

\compmeta{CounterfactualPlanningResult}{include/hacdcpf/analysis/counterfactual\_planning.hpp}
主要字段：\fld{schema}、\fld{baseline}（\texttt{CounterfactualMetricSet}）、
\fld{measures}、\fld{interactions}、\fld{methodology}；度量向量为经济运行成本、碳、可靠性
EENS、弹性加权 ENS 四维。

\section{核心数学模型}
\subsection{设备级承载力（DL/T 2041-2025）}
对每台变压器的下游供电区（BFS 展开），分布式电源可接入容量 $S_d$ 由代数公式
\begin{equation}
S_d = \frac{P - P_G + \beta\,S\cos\theta + P_{\text{ESS}} + \Delta P_{\text{ESS}}}{\tau_{\max}},
\end{equation}
其中 $P$ 为下游负荷，$P_G$ 为既有电源出力，$\beta$（\fld{single\_transformer\_beta}）为
单变压器可用容量系数，$S\cos\theta$ 为按功率因数折算的视在容量，$P_{\text{ESS}}$ 与
$\Delta P_{\text{ESS}}$ 为储能贡献，$\tau_{\max}$（\fld{n1\_loading\_limit}）为 N-1 负载率上限。
当 \fld{enable\_verification} 打开时，规范 §12 校核复用潮流、短路与谐波引擎（谐波为尽力而为）。

\subsection{多维薄弱环节筛查}
对实体 $i$、维度 $k=1,\dots,4$，先做阈值归一化，再取 Chebyshev 瓶颈严重度、最大最小共识度
与分歧度：
\begin{align}
p_{ik} &= \frac{\max(x_{ik},\,0)}{\tau_k}, \\
\sigma_i &= \max_k p_{ik}, \qquad
c_i = \min_k p_{ik}, \qquad
d_i = \max_k p_{ik} - \min_k p_{ik},
\end{align}
其中 $\tau_k$ 由 \fld{high\_pressure\_threshold} 等阈值给定，$\sigma_i$ 为瓶颈严重度、
$c_i$ 为共识度、$d_i$ 为维度分歧（range）。在此之上叠加 Pareto 偏序（\fld{pareto\_count}）
与时间维度风险投影（\fld{critical\_periods}）。该筛查为非补偿式：任一维度不达标不能被其他
维度补偿。

\subsection{反事实规划措施评估}
生成候选措施集合后，对每条措施 $j$ 施加到系统副本并重算四维度量向量
$\mathbf{m} = (C_{\text{op}},\, E_{\text{CO}_2},\, \text{EENS},\, \text{ENS})^\top$，收益为相对基线的差
\begin{equation}
\Delta\mathbf{m}_j = \mathbf{m}_{\text{base}} - \mathbf{m}_j,
\end{equation}
两两协同项由 $\Delta\mathbf{m}_{jl} - (\Delta\mathbf{m}_j + \Delta\mathbf{m}_l)$ 度量并写入
\fld{interactions}。经济/可靠性/弹性度量分别使用 \fld{energy\_price\_per\_mwh}、
\fld{reliability\_physical\_budget}、\fld{resilience\_horizon\_hours} 给定的有界口径。

\section{接口与参数}
\begin{paramtable}
\fld{default\_power\_factor} & $\cos\theta$ & — & $0\sim1$ & $0.95$ & 承载力默认功率因数\\
\fld{single\_transformer\_beta} & $\beta$ & — & $0\sim1$ & $0.80$ & 单变压器可用容量系数\\
\fld{n1\_loading\_limit} & $\tau_{\max}$ & — & $0\sim1$ & $1.0$ & N-1 负载率上限\\
\fld{enable\_verification} & — & 布尔 & — & \texttt{false} & 是否启用 §12 引擎校核\\
\fld{thd\_limit\_pct} & — & \% & — & $5.0$ & 谐波总畸变（THD）限值\\
\fld{top\_k} & $k$ & — & 整数 & $15$ & 输出前 $k$ 个薄弱环节\\
\fld{high\_pressure\_threshold} & $\tau$ & — & $0\sim1$ & $0.75$ & 高压力维度判定阈值\\
\fld{compatibility\_threshold} & — & — & $0\sim1$ & $0.50$ & 措施兼容判定阈值\\
\fld{conflict\_gap\_threshold} & — & — & $0\sim1$ & $0.60$ & 冲突间隙判定阈值\\
\fld{max\_measures} & — & — & 整数 & $5$ & 反事实候选措施上限\\
\fld{reliability\_physical\_budget} & — & — & 整数 & $40$ & 可靠性评估物理预算\\
\fld{resilience\_horizon\_hours} & — & h & — & $12$ & 弹性评估时窗\\
\fld{energy\_price\_per\_mwh} & — & \$/MWh & — & $100$ & 失负荷/电量折价\\
\end{paramtable}
反事实措施开关 \fld{enable\_line\_capacity}、\fld{enable\_storage}、\fld{enable\_tie\_switch}、
\fld{enable\_automation}、\fld{enable\_der} 控制候选措施族的启用；薄弱环节 \fld{mode} 取
\texttt{Planning}/\texttt{Operation}，\fld{target\_pressure} 为四维目标压力向量。

\section{验证与边界局限}
\subsection{验证与校核}
注册测试：\srcpath{test\_hosting\_capacity}、\srcpath{test\_multidimensional\_weak\_link}、
\srcpath{test\_counterfactual\_planning}，以及浏览器端 E2E
\srcpath{multidimensional\_weak\_link\_gui\_test}，覆盖承载力公式、四维筛查排序与反事实措施
评估的一致性。

\subsection{边界与局限（如实标注）}
\begin{itemize}
  \item 承载力：规范 §5/§6/§9 的系统级生产模拟、县域分解与系统—设备协调\textbf{未}实现；
        \fld{enable\_verification} 默认关闭，即默认仅规划筛查；谐波潮流校核为尽力而为；
  \item 薄弱环节：仅做筛查与决策路由，\textbf{不}宣称归因排放/停电量等于反事实投资收益；
        缺失维度不被当作零；
  \item 反事实：措施施加于系统副本，经济/可靠性/弹性指标使用选项给定的有界预算与时窗，
        非全时序生产模拟；
  \item 反事实规划暂无独立契约文档，口径以源码与本手册为准。
\end{itemize}
\end{document}
